% ============================================================================
% AI SECTION -- draft to be merged into the main paper as a single section.
% No abstract and no title: the \section below is the section itself.
% Style: anton_style.md. Items needing author confirmation are marked %% TODO.
% ============================================================================


\hypersetup{
  colorlinks=true,
  linkcolor=blue!55!black,
  urlcolor=blue!55!black,
  citecolor=blue!55!black
}

\setlength{\parindent}{0pt}
\setlength{\parskip}{6pt}
\setlist[itemize]{leftmargin=1.6em,itemsep=2pt,topsep=3pt}
\setlist[enumerate]{leftmargin=1.6em,itemsep=2pt,topsep=3pt}

\section{AI methodology and analysis}
\label{sec:ai}

{\color{red}AX: Overall, the opening before Subsection 11.1 is too long. It suffices to just have 1-2 paragraphs at most.}

The results in this paper were obtained in collaboration with an autonomous
AI research system. The cubic--quintic
construction and its certificate (Theorem~\ref{thm:cubic-quintic-upper-bound})
predate the current system and was found through conversation with GPT-5.5-Pro. The lower bound $\KG\ge6\pi/11$ was
discovered and first proved by the system; we have since verified the
argument ourselves, and the proof in this paper is our revision of the
system's write-up. The system also produced further improvements on both
sides of the problem, which we report in Section~\ref{sec:ai-results} but do
not state as theorems, because their certificates have not yet been human-verified. This section documents the methodology
(Section~\ref{sec:ai-methodology}), quantifies the run and its outcomes
(Section~\ref{sec:ai-results}), reconstructs how the lower bound was found
(Section~\ref{sec:ai-case-study}), and discusses what the record shows about
the strengths and weaknesses of current AI systems as mathematical
collaborators (Sections~\ref{sec:ai-discussion} and~\ref{sec:ai-future}).
{\color{red}AX: The discussion of cut-off dates belongs in some split of the methodology/discussions.
If you had to keep one paragraph, it would be this.
}


Recent successes of AI in mathematics, including on open problems, are of
a different kind from the work reported here. Evaluator-guided search has
produced record constructions~\cite{funsearch,novikov2025alphaevolvecodingagentscientific, charton2024patternboostconstructionsmathematicslittle},
and formal provers have shown proficiency at solving olympiad-level problems \cite{novikov2025alphaevolvecodingagentscientific}. In the
months surrounding our run, frontier models went further: an internal reasoning
model at OpenAI has disproved the Erd\H{o}s unit-distance conjecture
\cite{openai2026planarpointsets}, a mathematician has found a counterexample to the Jacobian conjecture using Claude Fable 5
\cite{jacobian}, and OpenAI has announced solutions, with formally
checked proofs, to ten open problems in mathematics and theoretical
computer science \cite{openai2026ten}. Most of these results involved resolving a single well-posed statement by construction. Moreover, none explicitly required significant long horizon reasoning. Finally, all of the recent results with frontier models were found fully autonomously, without sustained human intervention.
{\color{red}AX: Move to discussion. Also I'm not a fan of the ``in the months surrounding our run'' --- I think it's too fancy here. Let's just make this more standard ``recent and concurrent work'' that's happened to us.}

Our work is a more realistic example of how humans and AI can collaborate on mathematical research. The problem is more open ended, requiring a research program where over the course of weeks intermediate results accumulate and directions are opened and abandoned. Moreover, humans were consistently steering the system and providing feedback. We hope that our experiences outlined in this section will be useful to future mathematicians and AI researchers alike.
{\color{red}AX: Also move to discussion. It's also risky to explicitly claim that this is a ``more realistic example'', it could be more like: ``Our work is an extensive case-study ...'' or something like that?}



\subsection{Methodology}
\label{sec:ai-methodology}

{\color{red}AX: This is too dense. Tell Fable to explicitly split into the following paragraph groups: (1) AI models and compute: here you describe what models you used and you can also remark on the cut-off dates so that future readers have context. You can also add 1-2 sentences describing the server setup that you used.
(2) Harness: you can describe the harness system here.
(3) Run configurations
(4) Human interaction.
}

{\color{red}AX: Continuing from above: the current structure is kind of this split. It's more or less slapping on an explicit paragraph tag. However, the writing should be tightened to focus more on its assigned topic.}


The system uses a harness that couples two frontier language
models in distinct roles. A reasoning model (OpenAI GPT-5.5-Pro at maximal
reasoning effort, later GPT-5.6-Sol) serves as the brain, where it selects
directions, develops arguments, specifies experiments, and audits the
resulting claims. A coding agent (Anthropic Claude Code, running Claude Opus
and later Claude Fable~5) serves as the executor, where it maintains the project
repository, implements and runs experiments, retrieves literature, and
drives the reasoning model through a scripted interface. The harness was also adjusted repeatedly over the
run in response to observed failures. We describe its final form, to the
extent needed to interpret the record.


Work is organized into \emph{sessions} of a few hours, of which two to
five run concurrently on mutually distinct directions. No model context survives from
one session to the next, and instead continuity is carried by a single summary file that serves as
working memory and into which finished work is merged, and append-only
transcripts and experiment logs. A session passes through fixed phases:
orientation from the problem statement and summary files, selection of a focus, a research loop in
which the reasoning model requests computations, literature, or further
reasoning time, and a closing audit that produces a calibrated summary in
which every claim is tagged as proven, numerically supported, conjectural,
or heuristic. We reserve the word theorem, here and in the rest
of the paper, for statements that we have additionally checked ourselves.

The run started from a substantial base of our prior work rather than from
the problem alone. By June 2026 we had worked on the constant for several
months, and the harness was initialized with the results: the cubic--quintic scheme, the coefficient formulas and certification machinery, our accumulated notes, and a
problem statement listing candidate directions for both bounds. In
particular, the search space the run explored is a generalization of teh cubic--quintic scheme, and the possibility of converting
obstructions into lower bounds was raised, as a suggestion, in the problem
statement. The system therefore operated downstream of the mathematical
framework it was given, a point quantified in
Section~\ref{sec:ai-case-study}.

A human operator (one of the authors) steered the run asynchronously through
text files read at the start of, and during, every session. Steering
consisted of priorities, targets, audits of the run's own claims, and
occasional mathematical pointers; about forty dated directives were issued
over the run. The operator contributed no proof steps through this channel.

\subsection{Results}
\label{sec:ai-results}

The run spanned June~16 to July~24, 2026, It comprised
roughly 240 research sessions. The reasoning model was called 2{,}091
times, consuming about 152 million tokens at an estimated \$5{,}400 in
API cost; the coding agent ran on a flat-rate subscription (Claude Max). Floating-point searches ran
on a four-GPU node, and every certified computation was redone in interval
arithmetic on CPU. Table~\ref{tab:run} collects these figures, and
Table~\ref{tab:timeline} gives the timeline of the bounds.

\begin{table}[t]
\centering
\small
\begin{tabular}{@{}ll@{}}
\toprule
Duration & June 16 -- July 24, 2026 (main phase through July 4)\\
Research sessions & $\approx240$\\
Reasoning model & GPT-5.5-Pro, later GPT-5.6-Sol; maximal effort\\
Usage and cost & 2{,}091 calls; $\approx152$M tokens; $\approx\$5{,}400$\\
Coding agent & Claude Code (Opus, later Fable 5); subscription\\
Human steering & one operator; $\approx40$ dated directives\\
\bottomrule
\end{tabular}
\caption{\textbf{The run at a glance.} Aggregates computed from the harness
telemetry. Dollar figures are list-price estimates; coding-agent usage was
only partially metered.
}
\label{tab:run}
\end{table}

\begin{table}[t]
\centering
\small
\begin{tabular}{@{}lp{0.76\textwidth}@{}}
\toprule
Sessions & Event\\
\midrule
1 & run begins from the interval $1.6769566742\le\KG\le1.7818666070$
(the upper end is the cubic--quintic bound of this paper)\\
18 & upper-bound search plateaus; the operator directs a pivot to the
lower bound\\
44 & numerical upper bound $1.7802243$ recorded (later withdrawn)\\
55--57 & affine reframe; $\KG\ge6\pi/11$ derived\\
58--64 & $6\pi/11$ proof machine-audited; first \LaTeX{} write-up\\
p92, p117 & lower bounds $27\pi/49\approx1.7311$ and
$51\pi/92\approx1.7415$ machine-verified\\
follow-up & the session-44 value found unsupported by internal audit and
withdrawn\\
follow-up & upper bound $1.781801841033$ machine-certified\\
follow-up & upper bound $1.7813319810625639$ machine-certified;
adversarial review passed\\
\bottomrule
\end{tabular}
\caption{\textbf{Timeline of the bounds produced during the run.}
Machine-verified means accepted by the verification protocol of
Section~\ref{sec:ai-methodology}. The last three events fall in the short
follow-up phases, whose sessions we do not cite individually. Of the
run's results, only $\KG\ge6\pi/11$ has additionally been verified by the
authors; it is the result stated as a theorem in this paper.}
\label{tab:timeline}
\end{table}

The run's principal outcome, and the one this paper states as a theorem,
is the lower bound $\KG\ge6\pi/11$, whose discovery
Section~\ref{sec:ai-case-study} reconstructs. Beyond it, the run produced
machine-certified upper bounds of $1.781801841033$ and then
$1.7813319810625639$, improving on the cubic--quintic value, together with
the two stronger lower bounds listed in Table~\ref{tab:timeline}. We
report these as machine-verified claims rather than theorems: each passed
the full internal protocol, but we have not yet verified their
certificates ourselves.
%% TODO: reconcile with the main paper's discussion of concurrent work
%% (e.g. Heilman's computer-assisted improvement of Krivine's bound,
%% arXiv:2606.00247).

\subsection{Case study: the lower bound}
\label{sec:ai-case-study}

The lower bound is the run's main discovery. We reconstruct it here in
outline, highlighting the three human interventions that shaped it, since
the weaknesses discussed in Section~\ref{sec:ai-discussion} refer back to
them; the full logs of the sessions involved will be released with the rest
of the record.
{\color{red}AX: You don't have to refer to to-be-released artifacts here.
Tell Fable to write this version as though we're *not* releasing additional artifacts.
If and when we release artifacts, we'll strategically insert a few footnotes where needed.
}

{\color{red}AX: As currently written, I think the current section will be hard to track for anyone *not* intimately involved in the paper.
Granted, this is hard because it is about exploring a lot of concepts ...
Is it possible to extract a general/high-level for a proof template/sketch/idea so that we can ``set the scene'' + provide context?
\\
Additionally, I think you could consider adding a brief ``Overview`` paragraph of sorts that someone *not* working on Grothendieck can easily skim and get the big picture. This could be a high-level summarization of the subsequent ``The plateau``, ``First intervention'', etc. paragraphs.
}

\paragraph{The plateau.}
The run began on the upper bound, searching the limiting-scheme space of our
preprint for schemes whose correlation function admits a large
inverse-majorant parameter $\gamma$, since an admissible $\gamma$ gives
$\KG\le\pi/(2\gamma)$. From the hyperplane value
$\gamma=\log(1+\sqrt2)\approx0.8814$ it reached $\gamma\approx0.882$, and
there it stalled: in sessions 12 through 17 it opened six {\color{red}mechanically
distinct escape routes}, and each failed against the same obstruction,
namely that cubic coefficient mass cannot be cancelled without also reducing
the linear coefficient. Each failure was recorded and tagged, and the run's
response was to begin a seventh variation of the same search.

\paragraph{First intervention: turn the failures into a theorem.}
The operator redirected the run in session 18. Instead of attempting another
escape, it was to synthesize the accumulated failures into a proof that no
scheme in the class does much better: a ceiling $\gamma\le\Gamma$ valid for
the full class, with mixtures and arbitrary dimension, converts through the
optimality theorem of Naor and Regev \cite{naor2014krivine} into the lower
bound $\KG\ge\pi/(2\Gamma)$. The intervention supplied no new mathematics:
the problem statement had suggested this dualization from the outset, and
the run had accumulated exactly the required evidence without making the
connection. Session 18 identified the Naor--Regev theorem as the correct
transfer, proved a first ingredient valid in every dimension, and located
the difficulty precisely: the ceilings it could prove covered too narrow a
class of schemes.

\paragraph{Progress in two dimensions.}
Within days the run had proved a cubic ceiling of strength
$\Gamma\approx0.8866$ for two-dimensional schemes, with the {\color{red}keystone (?)}
inequality certified in interval arithmetic, and had sharpened the
constant numerically to $\Gamma\approx0.8846$. The extension to arbitrary dimension and to mixtures, which the
duality requires, did not follow: fourteen sessions in the range 32--51
attacked it, each reducing the gap to a sub-lemma, proving the sub-lemma in
a restricted setting, and leaving the global assembly open, while the
restricted results accumulated in memory as apparent progress. Two useful
facts emerged alongside this loop: the run verified the normalization of the
Naor--Regev transfer against the original paper, and it computed that
improving the Davie--Reeds~\cite{Davie1984LowerBoundKG,Reeds1991LowerBoundKG} bound $1.6769566742$ requires a universal ceiling
only at $\Gamma\le0.93669$, far above the two-dimensional value, so a
general ceiling could be far from sharp and still give a new lower bound.

\paragraph{Second intervention: the general bound, through theory.}
The operator ended this loop with a second redirection, on two grounds:
refinements of the two-dimensional result were not the object of interest,
since only a ceiling for the full class dualizes; and the gap was
conceptual, so further implementation would not close it. The run was to
spend several sessions on theory alone. Session 55 mapped the dependency
structure of the preceding weeks and identified the recurring error,
restricted closures recorded as global progress. Session 56 then produced
the {\color{red}decisive (?)} reframe. Every failed transport had been nonlinear in
quantities that do not average when schemes are mixed; the Hermite
coefficients $(b_1,b_3)$ of the correlation function, by contrast, are
linear functionals of the scheme, so a constraint that is affine in
$(b_1,b_3)$ passes to mixtures and to coefficientwise limits with no further
argument. Choosing the constraint for transportability, and requiring
equality at the extremal hyperplane point $(b_1,b_3)=(1,\tfrac16)$, gives
the one-parameter family of affine constraints
\[
        b_3\;\ge\;(1+\lambda)\,b_1-\Bigl(\lambda+\tfrac56\Bigr),
        \qquad \lambda\ge\tfrac12 ,
\]
where $\lambda$ indexes the slope of the line through the hyperplane point.
Each member forces a barrier $\Gamma_\lambda=(\lambda+\frac56)/(1+\lambda)$
on the admissible parameter, hence a bound $\KG\ge\pi/(2\Gamma_\lambda)$
conditional on proving that member, and every $\Gamma_\lambda$ below
$0.93669$ improves on Davie--Reeds.

\paragraph{Third intervention: one certified {\color{red}rung} before a better one.}
The tangent member $\lambda=\tfrac12$, worth $\KG\ge9\pi/16\approx1.767$,
requires a sharp chaos inequality that remains open, and once the family
existed the run's inclination was to work toward its strongest provable
member. The operator set the priority differently: what mattered was to
establish that the method yields any new lower bound at all, so the run was
to select a rung whose proof closes, certify it completely, and write it up,
with improvements inside the family treated as secondary. At $\lambda=1$ the
proof collapses to elementary structure: session 57 found that the
decomposition $h=(f+g)/2$, $k=(f-g)/2$ reduces the constraint, fiber by
one-dimensional fiber, to the ternary inequality~\eqref{eq:fiber-inequality},
%% TODO: fix cross-reference to the fiber inequality of the lower-bound section
whose closing case analysis is short, giving $\Gamma_1=\tfrac{11}{12}$ and
\[
        \KG\;\ge\;\frac{6\pi}{11}\;\approx\;1.7136 .
\]
Sessions 58--64 then subjected the chain to adversarial review, which
found and repaired the one substantive gap (the passage from finitely many
schemes to measurable mixtures, closed by a weak-$*$ lemma), a Jensen
inequality applied in the wrong direction, and a small defect in the
interval certificate's coverage, after which the run produced the first
complete write-up, about a day after the reframe. The stronger rungs
$27\pi/49$ and $51\pi/92$ of Table~\ref{tab:timeline} followed within a week
by the same method at smaller values of $\lambda$.



\paragraph{Attribution.}
The operator contributed no proof step anywhere in the chain. The operator's
contributions were the starting context of
Section~\ref{sec:ai-methodology} and the three interventions above, each a
redirection of effort rather than a piece of mathematics. The system
contributed the mathematics: the identification and verification of the
transfer theorem, the headroom computation, the two-dimensional ceiling, the
observation that affine constraints survive mixing, the family of lines, the
fiber reduction, the closure lemma, the certificates, and the initial
write-up. One further observation matters for what follows: nothing in the
session-56 reframe depended on information unavailable at the pivot a week
earlier, since the affineness of coefficients under mixing is elementary and
the headroom computation is one line. The binding constraint appears to have
been the system's disposition to continue its current program, not a missing
capability.
{\color{red} Name this something like ``human involvement'' or somehting along those lines. Also move the operator contributions first, and describe why they were necessary. Only then do you state the AI involvement.}


\subsection{Discussion}
\label{sec:ai-discussion}

The record supports a consistent capability profile: strong technical
execution, with persistent failures in the judgement that directs it. We
summarize the strengths, then the weaknesses, stated in the terms that the
interventions of Section~\ref{sec:ai-case-study} make concrete, and close
with what the record shows about the division of labor.

\subsubsection{Strengths}

Three strengths carried the run. The first is technical range with effective
recall: given a one-word pointer to Naor's work, the system identified the
relevant theorem, verified its normalization against the paper, and computed
the slack available for a lossy transfer; given a plateau, it produced six
structurally distinct escape mechanisms in two days. The second is local
invention: the affine family, the fiber reduction, the closure lemma, and
the certificate architecture were machine products, several arriving within
hours once the question was well posed. The third is speed and reliability
of execution: from the reframe to a machine-audited, interval-certified
proof took about a day, and most individual technical errors were caught by
the system's own verification passes. In one representative episode
(sessions p44 and p49), a claimed bound of $\KG\ge1.7754518$
was refuted within five hours by an adversarial session that computed its
true margin to be approximately $-8\times10^{-8}$, and the record was
corrected to the defensible value.

\subsubsection{Weaknesses}

Each intervention of Section~\ref{sec:ai-case-study} was needed because a
specific capability was absent. The paragraphs below name these
capabilities and add further instances from the record.

{\color{red}
AX: Name these paragraphs more explicitly to reflect a negative,
e.g., "Poor judgment of what is interesting", "Lack of high-level research skills", etc.}

\paragraph{Judgement of interestingness.}
The system did not reliably rank work by mathematical interest. The second
and third interventions are the clearest cases: the run treated refinements
of the two-dimensional ceiling as progress when only the general bound
mattered, and, once the family of rungs existed, it valued improvements
within the family over establishing the first certified rung. The same
tendency appeared at smaller scales throughout, from sessions spent on
$10^{-5}$-sized improvements to mechanisms the run already understood, down
to reported values tuned to the edges of their enclosures. The implication
for collaboration is that the objective of the research cannot yet be
delegated. The operator's judgements of interest were few, but each
redirected weeks of machine effort, and between them the system optimized,
competently, whatever was in front of it. On a problem of this size the
space of possible intermediate constructions is unbounded, and a ranking by
expected interest is what makes the search tractable; supplying that
ranking was the human's principal ongoing contribution.

\paragraph{Research skills.}
The system managed individual arguments well and the research process
poorly: it did not decide when to persist, when to switch, or what a
sequence of failures implied. The first intervention is the sharpest case,
since the dualization of failures into a lower bound was suggested in the
problem statement the system re-read at every session, the six failures
were in its memory, and the step was taken only when ordered. The record
contains three further episodes of the same shape, of fourteen, roughly two
dozen, and about a dozen sessions respectively (the two-dimensional
extension, near-sharp variants of the lower bound, and the July
certification effort), in each of which the run circled a target without
progress until the operator ended the loop. The implication is that
autonomous running time does not by itself convert into progress. Once the
accessible directions were exhausted, additional sessions multiplied
execution, not judgement, and the run's rate of real progress was set by
the frequency of strategic correction; additional compute would not have
changed it. Oversight of such systems is accordingly most valuable at the
level of strategy, which is precisely the level at which the system did not
correct itself.

\paragraph{Memory.}
The harness's file-based memory preserved statements more reliably than
their qualifications. The withdrawn upper bound of Table~\ref{tab:timeline}
is the clearest case. In session 44 the run recorded a numerical upper
bound of $1.7802243$, with caveats attached; over the nineteen days the
value stood in working memory, re-summarization detached the caveats from
the headline number, and a note defending the value was written three days
before the internal audit that refuted it, finding the score inflated by a
finite truncation window and the underlying correlation profile
infeasible. A related failure occurred inside verification, where a
dedicated replay of the $51\pi/92$ certificate (session p95) passed it
because the replay inherited the framing of the argument it was checking;
the invalidating error was found the same day by a session that rebuilt the
argument from scratch (session p102). We attribute both failures to scale.
The run produced hundreds of sessions and tens of millions of tokens of
intermediate mathematics, and each new session could be shown only a small
distillate of it. What to retain, at what resolution, and with which
qualifications attached was itself decided by the models, through repeated
summarization, and repeated summarization is exactly where caveats detach
and inherited framings go unquestioned. Unlike the previous two weaknesses,
this one looks amenable to engineering: memory for such systems plausibly
needs more structure than prose summaries provide, for instance a graph of
claims carrying provenance and qualifications, so that important
information remains retrievable without being repeatedly re-serialized
through a model's context window.

\paragraph{Origins of the weaknesses.}
We believe these gaps have a common origin in how current systems are
produced. Models are trained and evaluated on problems selected for being
small enough to pose, sandbox, and score, and nothing in that regime
exercises the allocation of weeks of effort across an open-ended search.
The corpus they learn mathematics from is, moreover, made of finished
mathematics: the literature records definitions, theorems, and polished
proofs, while the process that produced them, the dead ends, the decisions
to abandon or reframe, the judgement of what was worth attempting, is
largely absent from the written record. A system trained this way can be
expected to acquire technique and the form of finished mathematics while
remaining weak at exactly the process skills described above, and that is
what we observe.

\subsubsection{Division of labor}

{\color{red}Name this subsection like human operator involvement or something, to emphasize the human involvement first.}

A useful summary of the collaboration is obtained by counting who caught
which errors. Errors in individual claims, such as wrong constants, broken
certificates, and invalid steps, were caught almost entirely by the system
itself. Failure patterns spanning many sessions, such as the three circling
episodes and the detachment of caveats from claims, were caught by the
operator in every recorded instance and by the system in none. The system
reliably falsifies a claim once the claim is specified; deciding which
claims to distrust remained a human function throughout. Consistent with
this, the operator's leverage came less from mathematical suggestions, most
of which the system correctly tested and discarded, than from calibration:
setting targets, ending loops, and auditing the run's own reporting. We also
note a regularity with design consequences: every overclaiming incident in
the record sits at a target with nearly zero margin, whereas the results
that survived verification all carry slack. This suggests that systems of
this kind are most productive when aimed at robust targets, with sharp
constants attempted only from a secured base.

\subsection{Future directions}
\label{sec:ai-future}

\paragraph{Mathematics.}
The immediate items are the machine-verified claims of
Table~\ref{tab:timeline}: human verification of the two certified upper
bounds and of the lower bounds $27\pi/49$ and $51\pi/92$ would upgrade each
to a theorem. Beyond them, the run's internal analysis identifies the
tangent case $\lambda=\tfrac12$ of the affine family, equivalent to a sharp
degree-three chaos inequality, as the natural boundary of the method and
worth $\KG\ge9\pi/16$ if resolved.
{\color{red}AX: If you're brave, conjecture that \(K_G = \sqrt{\pi}\)}

\paragraph{Toward AI systems for long-horizon research.}
If the diagnosis of Section~\ref{sec:ai-discussion} is right, the
requirements for progress follow from it. First, mathematical
interestingness needs to become an object of study in its own right: the
judgement that selects which questions, intermediate constructions, and
targets are worth pursuing was the most consequential input the human
supplied, and there is at present no training signal or benchmark that
rewards it. Second, the process of research mathematics needs data. What
mathematicians actually do between problem statement and finished paper,
how they allocate effort, learn from repeated failure, and decide to
reframe, is barely recorded anywhere; records of long runs such as ours,
which preserve every step including the unsuccessful ones, are one way to
begin creating such data, and we release ours in that spirit. Third,
evaluation must move to horizons where these capabilities matter, treating
a weeks-long research program rather than a single problem as the unit of
assessment, with memory and calibration measured alongside correctness.
These conclusions carry the usual caution that they derive from one run on
one problem; replications on other open problems, successful or not, would
be the most useful follow-up work.

\paragraph{Human--AI collaboration.}
The run also suggests where collaboration of this kind is heading. Within
it, the contributions that required a human were not technical: the system
supplied the arguments, the computations, the certificates, and the first
write-up, while the human supplied the problem, the taste, and the
strategic corrections. If the technical abilities of these systems continue
to improve, this division will sharpen, and the mathematician's role in
such collaborations will concentrate into choosing what to work on, judging
which objects and targets matter, and deciding what to trust: directing
research that one does not personally execute. A similar shift is already
under way in software engineering. We expect working modes of this kind to
become common in mathematics, and the run documented here, with its
weaknesses as much as its results, is an early picture of what they will
require of the mathematician.

\subsection*{Statement on AI use}
{\color{red}AX: Any specific reason for an entire subsection instead of just a paragraph?}

Large language models were used extensively in the research reported in
this paper, in the roles described in this section: GPT-5.5-Pro and
GPT-5.6-Sol (OpenAI) as the reasoning model, and Claude Opus and Claude
Fable~5 (Anthropic) as the coding agent. AI assistance was also used in
preparing parts of this manuscript. All results stated as theorems have been
verified by the authors independently of these systems, and the account of
the run given in this section reflects the authors' own reading of the
record. The authors take full responsibility for the contents of this
paper.

